% RESEARCH_DRAFT_SHA256: 4f5c7e2fb6a55197c1e41ab7911386813ada5b75c1398e3b544eaf1cf118176d
% Layout and source-reference wrappers expanded for the web renderer.
\documentclass[11pt]{article}
\usepackage[margin=25mm]{geometry}
\usepackage{amsthm}
\usepackage{fontspec,unicode-math}
\setmainfont{Latin Modern Roman}
\setmathfont{Latin Modern Math}
\usepackage{amsmath,amssymb,mathtools,xurl,hyperref}
\hypersetup{colorlinks=true,urlcolor=blue,linkcolor=blue}
\setcounter{secnumdepth}{0}
\setlength{\parindent}{0pt}
\setlength{\parskip}{5pt}
\setlength{\emergencystretch}{3em}
\newcommand{\sref}[2]{\hyperlink{source:#1:#2}{[S#2]}}
\newcommand{\cataloguescope}[1]{\paragraph{Source record.}#1}
\newcommand{\floor}[1]{\left\lfloor#1\right\rfloor}
\newtheorem{proposition}{Proposition}
\newtheorem{lemma}[proposition]{Lemma}
\begin{document}
\paragraph{Research manuscript.} OpenAI \texttt{gpt-6-astra}, reasoning \texttt{ultra}; 11 October 2026.
\paragraph{Preserved source body.} The complete catalogue mathematical text and sources follow with layout and source links expanded, including source qualifications and its historical preparation notes. New research begins at the section headed Research attempt.

% UnsolvedMath ID 30006673; source catalogue code OWR-14299916-011
\section{Duplicate Unified-Theory Question}\label{30006673}
{\small UnsolvedMath ID 30006673 · Computer Science}
\subsection{Definitions and mathematical statement}

\paragraph{Shared notation.}Standard number systems are $\mathbb N,\mathbb Z,\mathbb Q,\mathbb R,\mathbb C$. Specialized notation and index ranges follow the cited source and the definitions below.


A learning theory may combine the geometry of parameter and representation spaces, probabilistic data laws, computational resource bounds, and training dynamics. A unified theory would need a specified common model and conclusions.
\paragraph{Record scope and qualifications.}
The catalogue title is retained as metadata. Named certificates, contracts, bounds or frameworks in that title are not asserted as results of the original source extract.
The supplied registry flags this extraction as unrecoverable. It is displayed as a literal source fragment; the definitions provide context rather than a replacement formal statement.

Source-record limitation: Source flags the statement as unrecoverable: This concluding rhetorical question repeats the unified-theory problem retained in OWR-14299916-002 and does not introduce a distinct mathematical target. No separate clean mathematical statement is supplied; this source ID preserves its provenance.

Related catalogue records: \href{https://mehmetmars7.github.io/Erdosproblems-llm-hunter/problem.html?type=open_problems&id=30006664}{30006664}. Separate source IDs are preserved for their provenance.
The source text below is rendered literally where necessary; damaged formulas and composite extracts are not silently repaired into a new mathematical claim.
The supplied catalogue marks this record as unrecoverable. Its extract and terminology are documented, but a complete original problem has not been reconstructed.
\paragraph{Statement recorded in the supplied source.}
The workshop demonstrated both the remarkable progress that has been achieved in understanding modern learning systems and the sub- stantial challenges that remain open. Whether machine learning will ultimately admit such a unified mathematical description remains an open question. Modern and Emerging Phenomena in Machine Learning 731 \href{https://doi.org/10.4171/owr/2026/11}{[S2]}
\paragraph{Source status.}
The supplied catalogue labels this source record open; that metadata was checked against the supplied snapshot on 2026-10-11. The label is preserved for provenance and is not an independent verification of present open status or mathematical correctness. The catalogue literature assessment is dated 2026-08-22.
\subsection{Short English statement}
The supplied fragment concerns duplicate Unified-Theory Question. Its complete mathematical target is not recovered.
\subsection{Sources}
\begin{itemize}
\item[S1] Ulam AI and the UnsolvedMath Contributors, UnsolvedMath: A Curated Collection of Open Mathematics Problems, record 30006673 (OWR-14299916-011). Catalogue curation and metadata: CC BY 4.0. The source extract is retained; definitions, English summaries and qualifications are editorial additions.
\url{https://huggingface.co/datasets/ulamai/UnsolvedMath}.
\item[S2] Modern and Emerging Phenomena in Machine Learning (2026), Oberwolfach Reports, DOI 10.4171/owr/2026/11, pp. 683--736
\url{https://doi.org/10.4171/owr/2026/11}.
\item Source-specific Oberwolfach Reports licences, including CC BY-SA 4.0 where applicable, are retained; see THIRD\_PARTY\_NOTICES.md.
\end{itemize}
\paragraph{Preparation and review.}
The original source and UnsolvedMath contributors retain attribution for the source extract. Codex (GPT-6; exact serving revision not exposed) assisted with definitions, English summaries, scope notes and formatting. No new human mathematical authorship or independently verified solution is asserted. Permanent source ID 30006673 and code OWR-14299916-011 checked against the complete supplied registry and current repository. Distinct upstream records, including repeated questions and broad or damaged extracts, are retained for complete coverage of the source CS category. A repeated record is not a new independent conjecture. Complete mathematical-equivalence and current-literature review has not been performed. Related canonical identities: 30006664. Preparation and final review follow contribution prompts C and D, with source gaps disclosed. Structural and rendering checks do not establish a new theorem or open status.
\label{end:30006673}

\section{Research attempt}
\subsection{Literature, source context and scope}
The full catalogue body, including qualifications and historical source status, is preserved above. The original workshop report was inspected in the locally cached publisher PDF text. This attempt distinguishes the source question from any explicitly restricted model introduced below. Repeated records retain their own source IDs and do not become independent conjectures. Elementary results and attributed literature results are not claimed as new original-target progress.

\subsection{A finite-observation route and an identifiability limit}
This concluding rhetorical fragment repeats the unified-theory question. As a limited test of a universal empirical-prediction route, consider binary labeling on a finite input space. Training reveals labels at a set $S$, but not at one further point $x_*$. The hypothesis class permits both labels at $x_*$ independently of the observed labels. A test law gives probability $\alpha\in(0,1]$ to $x_*$ and puts the remaining probability on $S$.
\begin{proposition}
For every possibly randomized learner using only the observed labels, some consistent target gives expected test error at least $\alpha/2$. This minimax bound is attained by predicting the observed labels on $S$ and a fair random bit at $x_*$.
\end{proposition}
\begin{proof}
Choose two targets identical on $S$ and opposite at $x_*$. Their training observations, and therefore the learner's output distribution, are identical. Let $p$ be the probability the learner predicts one at $x_*$. Its conditional error at that point is $p$ under the zero target and $1-p$ under the one target. Their maximum is at least $1/2$; multiplying by the test mass $\alpha$ gives the lower bound, with any additional errors on $S$ nonnegative. The declared fair-bit learner makes no errors on $S$ and has conditional error exactly $1/2$ at the new point under either target, attaining $\alpha/2$.
\end{proof}
Thus observations alone cannot identify arbitrary unseen behavior. A unifying theory can certainly include this uncertainty or impose structural assumptions linking unseen and seen cases; the calculation does not refute the possibility of such a theory. It blocks only the inference that finite agreement by itself supplies a universal determinate extrapolation.
\subsection{Verification and remaining gaps}
For $\alpha=1/4$, probabilities $p=0,1/4,1/2,3/4,1$ yield the predicted worst-case errors, minimized at $1/8$. The result is a standard elementary indistinguishability bound, not new original-target progress. The source does not define what would count as a coherent unified theory, and this limited model does not furnish one.

\section{References and source attribution}
The primary Oberwolfach workshop report, \url{https://doi.org/10.4171/owr/2026/11}, was inspected on 11 October 2026. Relevant source passages: 2026/11 pp. 730--731, concluding repeated unified-description question. Catalogue and contributor attribution remain in the original Sources section. Mathematical arguments labeled as this attempt are proved above; cited external results remain attributed to their authors and are not independently reproved or claimed as our discoveries.


\paragraph{Integration review.} Codex (GPT-6; exact serving revision not exposed) reviewed the written arguments and expanded layout and source-reference wrappers. Repository metadata and the production web renderer were checked using Python 3.13 and JavaScript with MathJax 3.2.2. This model-assisted review is not independent human proof verification; no PDF compilation or formal proof checking is asserted.

\section{Completion Estimate}
\noindent\textbf{COMPLETION: 0\%}
No substantiated new progress toward the entire original target is claimed. The worked results and failed route are delimited above; a known result, illustrative model, or auxiliary obstruction is not counted as original-target completion.

\section{Final status and attribution}
\textbf{Attempt status: unresolved.} This describes this research attempt, not a new ruling on the source catalogue's status or on external literature claims.
The mathematics and this self-audit were generated by OpenAI \texttt{gpt-6-astra}, reasoning \texttt{ultra}, on 11 October 2026. Codex (GPT-6; exact serving revision not exposed) orchestrated the batch and contributed formatting/repository review. No human mathematical authorship, independent human proof checking, or formal verification is claimed. The historical preparation paragraphs in the reproduced catalogue body concern that body only.

Contribution prompt A was applied to the complete definition and supplied source record, followed by the adversarial checks of prompt B. Local reproduction files are \texttt{batch011/checks.py}, \texttt{computations.json}, and the per-ID audit JSON. They were not public links at preparation time. Numerical checks are finite and do not prove unrestricted results. No LaTeX engine was found locally; compilation was not performed by this research agent.

\end{document}
